在命题 \ref{prop:pom-diagonal-rate-rank-one-refresh-decomposition} 的秩一刷新表示下，最优核的首达时间与电网络对象出现一维闭合。
其要点在于：每一步仅包含“状态依赖保持”与“全局刷新采样”两种机制，
从而首达时间分布可分解为几何次数的几何停留时间之和；同时，由可逆性诱导的电导完全图为秩一乘积，Kron 还原退化为星形网络。

\paragraph{首达时间的概率母函数与复合几何表示}
沿用命题 \ref{prop:pom-diagonal-rate-rank-one-refresh-decomposition} 的记号，
令 \((X_t)_{t\ge 0}\) 为以核 \(K_\delta\) 驱动、起点 \(X_0=x\) 的离散时间马尔可夫链。
对目标点 \(y\in X\) 定义首次到达时间
\[
\tau_y:=\inf\{t\ge 0:\ X_t=y\}.
\]
并对每个起点 \(x\in X\) 定义“首次刷新等待时间”
\[
H_x:=\inf\{t\ge 1:\ \text{在第 \(t\) 步触发刷新}\}.
\]

\begin{lemma}[几何停留时间的母函数]\label{lem:pom-diagonal-rate-refresh-holding-time-pgf}
\leanverified{paper\_pom\_diagonal\_rate\_refresh\_holding\_time\_pgf}
在秩一刷新表示
\(
K_\delta=\mathrm{diag}(1-r_\delta)+r_\delta\,\pi_\delta^{\mathsf T}
\)
下，对每个 \(x\in X\)，随机变量 \(H_x\) 服从参数为 \(r_\delta(x)\) 的几何分布（取值于 \(\{1,2,\dots\}\)），其概率母函数为
\[
G_x(s):=\EE[s^{H_x}]=\frac{r_\delta(x)s}{1-(1-r_\delta(x))s},
\qquad |s|<\frac{1}{1-r_\delta(x)}.
\]
\end{lemma}
\begin{proof}
由定义，链在每一步以概率 \(1-r_\delta(x)\) 保持在 \(x\)，以概率 \(r_\delta(x)\) 触发刷新；\(H_x\) 即“首次刷新”所需的步数。
因此 \(\mathbb P(H_x=t)=(1-r_\delta(x))^{t-1}r_\delta(x)\)（\(t\ge 1\)），求和得到母函数闭式。证毕。
\end{proof}

\begin{theorem}[首达时间的概率母函数闭式与复合几何分解]\label{thm:pom-diagonal-rate-refresh-hitting-time-pgf-closed}
对任意互异 \(x,y\in X\)，\(\tau_y\) 的概率母函数在 \(s=0\) 的邻域内解析，并满足严格闭式
\[
\EE_x[s^{\tau_y}]
=G_x(s)\cdot \frac{\pi_\delta(y)}{1-(1-\pi_\delta(y))\,\overline G_{-y}(s)},
\qquad
\overline G_{-y}(s):=\sum_{z\ne y}\frac{\pi_\delta(z)}{1-\pi_\delta(y)}\,G_z(s).
\]
等价地，\(\tau_y\) 具有如下复合分布表示：令 \(N\sim\mathrm{Geom}(\pi_\delta(y))\) 为取值于 \(\{0,1,2,\dots\}\) 的几何分布（失败轮数），并令
\(
Z_i\overset{\text{\rm i.i.d.}}{\sim}\pi_\delta(\cdot\mid\cdot\ne y)
\)
且与 \(N\) 及 \(\{H_z\}\) 相互独立，则
\[
\tau_y\ \overset{d}{=}\ H_x+\sum_{i=1}^{N} H_{Z_i}.
\]
\end{theorem}
\begin{proof}
由秩一刷新表示，若 \(X_t\ne y\)，则离开当前状态的唯一机制是触发一次刷新并从 \(\pi_\delta\) 采样下一状态。
从 \(x\ne y\) 出发，首次刷新前的步数为 \(H_x\)（引理 \ref{lem:pom-diagonal-rate-refresh-holding-time-pgf}）。
刷新后以概率 \(\pi_\delta(y)\) 直接命中 \(y\) 并终止；否则落在某个 \(z\ne y\)，并从该 \(z\) 重新开始同一过程。
因此失败轮数 \(N\) 服从 \(\mathrm{Geom}(\pi_\delta(y))\)，失败轮的落点序列 \(Z_i\) 为条件化背景分布 \(\pi_\delta(\cdot\mid\cdot\ne y)\) 的独立样本；
每轮在落点 \(z\) 的等待时间为 \(H_z\)，且与 \(N,Z_i\) 独立。
由复合独立性得到分布等式；对母函数取期望并使用 \(\EE[t^N]=\pi_\delta(y)/(1-(1-\pi_\delta(y))t)\) 得闭式。证毕。
\end{proof}

\begin{corollary}[均值命中时闭式]\label{cor:pom-diagonal-rate-refresh-hitting-time-mean-closed}
\leanverified{paper\_pom\_diagonal\_rate\_refresh\_hitting\_time\_mean\_closed}
对任意互异 \(x,y\in X\)，有
\[
\EE_x[\tau_y]=\frac{1}{r_\delta(x)}+\frac{1}{\pi_\delta(y)}\sum_{z\ne y}\frac{\pi_\delta(z)}{r_\delta(z)}.
\]
\end{corollary}
\begin{proof}
由定理 \ref{thm:pom-diagonal-rate-refresh-hitting-time-pgf-closed} 的复合分布表示，\(\EE[H_x]=1/r_\delta(x)\) 且
\(
\EE[N]=(1-\pi_\delta(y))/\pi_\delta(y)
\)。
并且
\(
\EE[H_{Z_1}]=\sum_{z\ne y}\frac{\pi_\delta(z)}{1-\pi_\delta(y)}\cdot \frac{1}{r_\delta(z)}
\)。
代回即得所示闭式。证毕。
\end{proof}

\paragraph{删点 Laguerre 多项式与几何谱分解}
当核进一步专化为最优核 \(K^\star_\delta\) 的显式参数化（命题 \ref{prop:pom-diagonal-rate-diagonal-closure-rankone}）时，
首达时间母函数可被压缩为一个删点 Laguerre 多项式的有理式；
其分母的实根交错把首达时间分布的全部信息压缩为一组落在 \((0,1)\) 的几何谱参数 \(\{\lambda_{y,i}\}\) 以及相应的残数系数。

\begin{theorem}[首达时间母函数的删点 Laguerre 闭式]\label{thm:pom-diagonal-rate-refresh-hitting-time-deleted-laguerre-pgf}
\leanverified{paper\_pom\_diagonal\_rate\_refresh\_hitting\_time\_deleted\_laguerre\_pgf}
在命题 \ref{prop:pom-diagonal-rate-diagonal-closure-rankone} 的设定下，固定互异 \(x,y\in X\)，并令 \((X_t)_{t\ge 0}\) 由最优核 \(K^\star_\delta\) 驱动且 \(X_0=x\)。
定义首达时间 \(\tau_y:=\inf\{t\ge 0:\ X_t=y\}\)。
令删点多项式
\[
P_{-y}(z):=\prod_{u\in X\setminus\{y\}}(t_u-z),\qquad
D_y(z):=zP_{-y}'(z)+\kappa_\delta P_{-y}(z),
\qquad
P_{-x,-y}(z):=\prod_{u\in X\setminus\{x,y\}}(t_u-z).
\]
则对 \(|s|\) 充分小有严格恒等式
\[
\EE_x[s^{\tau_y}]
=
\frac{\kappa_\delta\,s\,(t_x-\kappa_\delta)}{t_y-\kappa_\delta}\cdot
\frac{P_{-x,-y}(\kappa_\delta s)}{D_y(\kappa_\delta s)}.
\]
\end{theorem}
\begin{proof}
由命题 \ref{prop:pom-diagonal-rate-diagonal-closure-rankone}，最优核满足秩一刷新分解
\[
K^\star_\delta(z\mid x)=(1-r_\delta(x))\mathbf 1_{\{z=x\}}+r_\delta(x)\pi_\delta(z),
\qquad
r_\delta(x)=\frac{t_x-\kappa_\delta}{t_x},\quad
\pi_\delta(z)=\frac{1}{t_z-\kappa_\delta}.
\]
沿用定理 \ref{thm:pom-diagonal-rate-refresh-hitting-time-pgf-closed} 的记号并取 \(f_x(s):=\EE_x[s^{\tau_y}]\)（\(x\ne y\)），则由一步分解得到线性方程
\(
(I-sT^{(y)})f=sq^{(y)}
\)
（\(T^{(y)}\) 为删去 \(y\) 的次转移矩阵，\(q^{(y)}(x)=K^\star_\delta(y\mid x)\)）。
在秩一结构下可直接化为
\[
f_x(s)
=\frac{s\,r_\delta(x)\pi_\delta(y)}{1-s(1-r_\delta(x))}
\cdot
\frac{1}{1-s\sum_{z\ne y}\frac{\pi_\delta(z)r_\delta(z)}{1-s(1-r_\delta(z))}}.
\]
代入 \(\pi_\delta(z)r_\delta(z)=1/t_z\) 与 \(1-s(1-r_\delta(z))=(t_z-\kappa_\delta s)/t_z\) 得
\[
\sum_{z\ne y}\frac{\pi_\delta(z)r_\delta(z)}{1-s(1-r_\delta(z))}
=\sum_{z\ne y}\frac{1}{t_z-\kappa_\delta s}.
\]
令 \(z=\kappa_\delta s\) 并用对数导数恒等式
\(
\sum_{u\ne y}\frac{1}{t_u-z}=-\frac{P_{-y}'(z)}{P_{-y}(z)}
\)
化简得到
\[
\frac{1}{1-s\sum_{z\ne y}\frac{1}{t_z-\kappa_\delta s}}
=\frac{\kappa_\delta P_{-y}(\kappa_\delta s)}{D_y(\kappa_\delta s)}.
\]
再注意 \(P_{-y}(\kappa_\delta s)=(t_x-\kappa_\delta s)P_{-x,-y}(\kappa_\delta s)\) 并与前因子中的 \((t_x-\kappa_\delta s)^{-1}\) 消去，整理即得所示恒等式。证毕。
\end{proof}

\begin{theorem}[吸收 Laguerre 多项式的实根交错]\label{thm:pom-diagonal-rate-refresh-absorbing-laguerre-interlacing}
\leanverified{paper\_pom\_diagonal\_rate\_refresh\_absorbing\_laguerre\_interlacing}
在定理 \ref{thm:pom-diagonal-rate-refresh-hitting-time-deleted-laguerre-pgf} 的设定下，
设 \(\{t_u:\ u\ne y\}\) 的升序排列为 \(t_{(1)}<\cdots<t_{(n-1)}\)（其中 \(n:=|X|\ge 2\)）。
则 \(D_y\) 的全部零点均为单实根，记为
\(
\kappa_\delta<z_{y,1}<\cdots<z_{y,n-1}
\)
并满足严格交错
\[
\kappa_\delta<z_{y,1}<t_{(1)}<z_{y,2}<t_{(2)}<\cdots<z_{y,n-1}<t_{(n-1)}.
\]
特别地，\(\lambda_{y,i}:=\kappa_\delta/z_{y,i}\in(0,1)\)。
\end{theorem}
\begin{proof}
注意 \(D_y(z)=0\) 等价于世俗方程
\[
\sum_{u\ne y}\frac{1}{t_u-z}=\frac{\kappa_\delta}{z},
\]
其左端在每个开区间 \((0,t_{(1)})\) 与 \((t_{(k)},t_{(k+1)})\) 上连续严格递增并在端点处发散，
而右端在 \((0,\infty)\) 上严格递减且为正。
由介值定理可知在每个 \((t_{(k)},t_{(k+1)})\) 内存在唯一解，并在 \((0,t_{(1)})\) 内存在唯一解。
\par
为定位该最小解在 \((\kappa_\delta,t_{(1)})\) 内，利用命题 \ref{prop:pom-diagonal-rate-diagonal-closure-rankone} 的归一化条件
\(
\sum_{u\in X}(t_u-\kappa_\delta)^{-1}=1
\)
得到
\(
\sum_{u\ne y}(t_u-\kappa_\delta)^{-1}=1-(t_y-\kappa_\delta)^{-1}<1=\kappa_\delta/\kappa_\delta
\)。
结合左端在 \((0,t_{(1)})\) 上严格递增与 \(t_{(1)}>\kappa_\delta\) 可知该唯一解必落在 \((\kappa_\delta,t_{(1)})\)。
各区间解的唯一性同时推出全部根均为单根，从而得到所示交错与 \(\lambda_{y,i}\in(0,1)\)。证毕。
\end{proof}

\begin{theorem}[几何谱分解表示]\label{thm:pom-diagonal-rate-refresh-hitting-time-geometric-spectral-decomp}
在定理 \ref{thm:pom-diagonal-rate-refresh-absorbing-laguerre-interlacing} 的设定下，
对任意互异 \(x,y\in X\) 存在唯一系数 \(\{\omega_i^{(x,y)}\}_{i=1}^{n-1}\subset\RR\)（\(\sum_i\omega_i^{(x,y)}=1\)）使对一切 \(k\ge 1\) 成立
\[
\PP_x(\tau_y=k)=\sum_{i=1}^{n-1}\omega_i^{(x,y)}\,(1-\lambda_{y,i})\,\lambda_{y,i}^{k-1}.
\]
并且系数具有显式残数公式
\[
\omega_i^{(x,y)}
=
-\frac{t_x-\kappa_\delta}{t_y-\kappa_\delta}\cdot
\frac{\kappa_\delta\,P_{-x,-y}(z_{y,i})}{(z_{y,i}-\kappa_\delta)\,D_y'(z_{y,i})}.
\]
\leanverified{paper\_pom\_diagonal\_rate\_refresh\_hitting\_time\_geometric\_spectral\_decomp}
\end{theorem}
\begin{proof}
由定理 \ref{thm:pom-diagonal-rate-refresh-hitting-time-deleted-laguerre-pgf} 可知 \(\EE_x[s^{\tau_y}]/s\) 为关于 \(z=\kappa_\delta s\) 的有理函数
\[
F(z)=\frac{t_x-\kappa_\delta}{t_y-\kappa_\delta}\cdot \kappa_\delta\,\frac{P_{-x,-y}(z)}{D_y(z)}.
\]
由定理 \ref{thm:pom-diagonal-rate-refresh-absorbing-laguerre-interlacing}，分母 \(D_y\) 具有互异实根 \(\{z_{y,i}\}_{i=1}^{n-1}\)，且 \(\deg P_{-x,-y}=n-2< n-1=\deg D_y\)，故 \(F\) 可作部分分式分解并回代 \(z=\kappa_\delta s\) 得
\[
\EE_x[s^{\tau_y}]
=\sum_{i=1}^{n-1}c_i^{(x,y)}\,\frac{s}{1-\lambda_{y,i}s},
\qquad
c_i^{(x,y)}:=-\frac{t_x-\kappa_\delta}{t_y-\kappa_\delta}\cdot
\frac{\kappa_\delta\,P_{-x,-y}(z_{y,i})}{z_{y,i}D_y'(z_{y,i})}.
\]
将每一项改写为几何母函数 \(\frac{(1-\lambda)s}{1-\lambda s}=\sum_{k\ge 1}(1-\lambda)\lambda^{k-1}s^k\) 的线性组合即得所示表示，其中 \(\omega_i^{(x,y)}=c_i^{(x,y)}/(1-\lambda_{y,i})\)，即为残数公式。
归一化 \(\sum_i\omega_i^{(x,y)}=1\) 由 \(s\uparrow 1\) 时 \(\EE_x[s^{\tau_y}]\to 1\) 及每一几何母函数在 \(s=1\) 取值为 \(1\) 得出；
唯一性由 \(\{\lambda_{y,i}\}\) 互异而母函数表示的线性无关性给出。证毕。
\end{proof}

\begin{corollary}[尾概率的几何谱展开]\label{cor:pom-diagonal-rate-refresh-hitting-time-tail-spectral}
在定理 \ref{thm:pom-diagonal-rate-refresh-hitting-time-geometric-spectral-decomp} 的设定下，令尾概率 \(S_k^{(x,y)}:=\PP_x(\tau_y>k)\)。
则对一切 \(k\ge 0\) 有
\[
S_k^{(x,y)}=\sum_{i=1}^{n-1}\omega_i^{(x,y)}\,\lambda_{y,i}^{k}.
\]
\leanverified{paper\_pom\_diagonal\_rate\_refresh\_hitting\_time\_tail\_spectral}
\end{corollary}
\begin{proof}
由定理 \ref{thm:pom-diagonal-rate-refresh-hitting-time-geometric-spectral-decomp} 的表示逐项求和即可得到尾概率展开。证毕。
\end{proof}

\begin{remark}[系数正性与完全单调尾]\label{rem:pom-diagonal-rate-refresh-hitting-time-tail-hausdorff-condition}
在推论 \ref{cor:pom-diagonal-rate-refresh-hitting-time-tail-spectral} 的谱展开中，\(\lambda_{y,i}\in(0,1)\) 由实根交错强制；
然而系数 \(\omega_i^{(x,y)}\) 的符号并不由该交错性质单独决定。
若进一步满足 \(\omega_i^{(x,y)}\ge 0\)（从而 \(\sum_i\omega_i^{(x,y)}=1\) 给出 \([0,1)\) 上的离散概率测度 \(\sum_i\omega_i^{(x,y)}\delta_{\lambda_{y,i}}\)），则 \(\{S_k^{(x,y)}\}_{k\ge 0}\) 为 Hausdorff 矩序列并满足完全单调性
\[
(-1)^m(\Delta^m S^{(x,y)})_k\ge 0\qquad(m,k\ge 0),
\]
并且风险率 \(h_k^{(x,y)}:=\PP_x(\tau_y=k\mid \tau_y\ge k)\) 对 \(k\ge 1\) 具有递减性 \(h_{k+1}^{(x,y)}\le h_k^{(x,y)}\)。
\end{remark}

% Auto-split to keep each TeX source < 600 lines.

\paragraph{接受—刷新强平稳时刻与分离距离的谱闭式}\label{para:pom-diagonal-rate-refresh-strong-stationary-separation}
在最优核 \(K^\star_\delta\) 的“保持 \(+\) 全局刷新”秩一结构下（命题 \ref{prop:pom-diagonal-rate-diagonal-closure-rankone}），
可以在同一代数系统内同时封闭三类对象：
强平稳时刻、分离距离及其谱侧权重。
其核心事实是：沿刷新时刻做拒绝采样，可以把背景分布 \(\pi_\delta\) 精确重加权为平稳分布 \(w\)，并由一类 halting 状态把分离距离的最坏点固定在时间上。

\begin{definition}[接受—刷新强平稳时刻]\label{def:pom-diagonal-rate-accept-refresh-sst}
在命题 \ref{prop:pom-diagonal-rate-diagonal-closure-rankone} 的设定下，
令 \((X_t)_{t\ge 0}\) 为以最优核 \(K^\star_\delta\) 驱动、起点 \(X_0=x\) 的离散时间马尔可夫链。
由命题 \ref{prop:pom-diagonal-rate-diagonal-closure-rankone} 的秩一闭包，对任意 \(x,y\in X\) 有
\[
K^\star_\delta(y\mid x)
=r_\delta(x)\,\pi_\delta(y)+\bigl(1-r_\delta(x)\bigr)\mathbf 1_{\{y=x\}},
\qquad
r_\delta(x)=\frac{t_x-\kappa_\delta}{t_x},
\quad
\pi_\delta(y)=\frac{1}{t_y-\kappa_\delta},
\]
并满足归一化 \(\sum_{y\in X}\pi_\delta(y)=1\)。
\par
记最小刷新率与其最小点集为
\[
c_\delta:=\min_{z\in X} r_\delta(z)\in(0,1),
\qquad
\mathcal H_\delta:=\arg\min_{z\in X} r_\delta(z)=\arg\min_{z\in X} t_z.
\]
取定一个 \(h\in\mathcal H_\delta\) 并令接受率
\[
a_\delta(y):=\frac{c_\delta}{r_\delta(y)}\in(0,1]\qquad(y\in X),
\qquad a_\delta(h)=1.
\]
在每一次“触发刷新”的步上，链产生新状态 \(Y:=X_t\sim\pi_\delta\)（条件于刷新事件），并独立地以概率 \(a_\delta(Y)\) 停止。
令 \(T\) 为首次停止时刻（即首次被接受的刷新步）。
\end{definition}

\begin{theorem}[强平稳性]\label{thm:pom-diagonal-rate-accept-refresh-sst-strong}
在定义 \ref{def:pom-diagonal-rate-accept-refresh-sst} 的设定下，
\(T\) 是 \(K^\star_\delta\) 的强平稳时刻：对每个起点 \(x\in X\)，随机变量 \(X_T\) 服从平稳分布 \(w\)，并且与 \(T\) 独立。
\end{theorem}
\begin{proof}
由命题 \ref{prop:pom-diagonal-rate-diagonal-closure-rankone} 的显式参数化，
对任意 \(y\in X\) 有
\(
w(y)\propto \dfrac{t_y}{(t_y-\kappa_\delta)^2}
\)
（见定理 \ref{thm:pom-diagonal-rate-diagonal-statistics-complete} 的证明要点）。
另一方面，
\[
\pi_\delta(y)\,\frac{1}{r_\delta(y)}
=\frac{1}{t_y-\kappa_\delta}\cdot\frac{t_y}{t_y-\kappa_\delta}
=\frac{t_y}{(t_y-\kappa_\delta)^2},
\]
从而 \(w\) 与向量 \(\bigl(\pi_\delta(y)/r_\delta(y)\bigr)_{y\in X}\) 同比例。
因此在任意一次刷新中，接受后的状态分布满足
\[
\PP(Y=y\mid \text{\rm 刷新且被接受})
\propto \pi_\delta(y)\,a_\delta(y)
=c_\delta\,\frac{\pi_\delta(y)}{r_\delta(y)}
\propto w(y),
\]
故 \({\rm Law}(Y\mid \text{\rm 刷新且被接受})=w\)，且该分布与刷新发生的时间无关。
由于停止只可能发生在刷新步，且“是否停止”仅依赖于该步的刷新样本 \(Y\) 与独立抛硬币，
得 \(\PP_x(T=t,X_T=y)=\PP_x(T=t)\,w(y)\)。
等价地，\(X_T\sim w\) 且与 \(T\) 独立。证毕。
\end{proof}
\leanverified{paper\_pom\_diagonal\_rate\_accept\_refresh\_sst\_strong}

\begin{proposition}[halting 状态]\label{prop:pom-diagonal-rate-accept-refresh-halting-state}
在定义 \ref{def:pom-diagonal-rate-accept-refresh-sst} 的设定下，若起点 \(x\ne h\)，则对任意 \(m\ge 0\) 有事件包含
\[
\{X_m=h\}\subseteq \{T\le m\}.
\]
\end{proposition}
\begin{proof}
当 \(x\ne h\) 时，要在时刻 \(m\) 处于 \(h\) 必须在某个时刻 \(t\in\{1,\dots,m\}\) 发生刷新并采样到 \(h\)（非刷新步只能自环）。
而一旦刷新样本为 \(h\)，由 \(a_\delta(h)=1\) 得该刷新步必被接受并立即停止，故 \(T\le t\le m\)。证毕。
\end{proof}
\leanverified{paper\_pom\_diagonal\_rate\_accept\_refresh\_halting\_state}

\begin{theorem}[精确分离距离与时间一致的最坏点]\label{thm:pom-diagonal-rate-accept-refresh-separation-exact}
\leanverified{paper\_pom\_diagonal\_rate\_accept\_refresh\_separation\_exact}
在定义 \ref{def:pom-diagonal-rate-accept-refresh-sst} 的设定下，
对任意起点 \(x\in X\setminus\{h\}\) 与任意 \(m\ge 0\)，分离距离
\[
\mathrm{sep}_x(m):=\max_{y\in X}\Bigl(1-\frac{(K^\star_\delta)^m(y\mid x)}{w(y)}\Bigr)
\]
满足严格恒等式
\[
\mathrm{sep}_x(m)=\PP_x(T>m)=1-\frac{(K^\star_\delta)^m(h\mid x)}{w(h)}.
\]
并且最坏点在时间上不变：对所有 \(m\ge 0\)，有
\[
h\in\arg\min_{y\in X}\frac{(K^\star_\delta)^m(y\mid x)}{w(y)}.
\]
\leanverified{paper\_pom\_diagonal\_rate\_accept\_refresh\_separation\_exact}
\end{theorem}
\begin{proof}
由定理 \ref{thm:pom-diagonal-rate-accept-refresh-sst-strong} 的强平稳性，
\(
\PP_x(T\le m,\,X_m=y)=\PP_x(T\le m)\,w(y)
\)
对一切 \(y\in X\) 成立。
因此
\[
\frac{(K^\star_\delta)^m(y\mid x)}{w(y)}
=\frac{\PP_x(X_m=y)}{w(y)}
\ge \PP_x(T\le m),
\]
从而 \(\mathrm{sep}_x(m)\le \PP_x(T>m)\)。
\par
另一方面，由命题 \ref{prop:pom-diagonal-rate-accept-refresh-halting-state}，对 \(x\ne h\) 有
\(\PP_x(X_m=h,\,T>m)=0\)，故
\[
(K^\star_\delta)^m(h\mid x)
=\PP_x(T\le m,\,X_m=h)
=\PP_x(T\le m)\,w(h),
\]
即 \((K^\star_\delta)^m(h\mid x)/w(h)=\PP_x(T\le m)\)。
代回即得
\(
1-(K^\star_\delta)^m(h\mid x)/w(h)=\PP_x(T>m)
\)
并与上界合并得到所示恒等式。
最后，由 \((K^\star_\delta)^m(y\mid x)/w(y)\ge \PP_x(T\le m)=(K^\star_\delta)^m(h\mid x)/w(h)\)，知 \(h\) 总实现最小比值，故为时间一致的最坏点之一。证毕。
\end{proof}

\begin{corollary}[严格分离]\label{cor:pom-diagonal-rate-accept-refresh-separation-strict}
\leanverified{paper\_pom\_diagonal\_rate\_accept\_refresh\_separation\_strict}
在定理 \ref{thm:pom-diagonal-rate-accept-refresh-separation-exact} 的设定下，
若 \(\mathcal H_\delta=\{h\}\) 为单点（等价地，\(t_h< t_y\) 对一切 \(y\ne h\) 成立），则对任意 \(x\ne h\)、任意 \(m\ge 0\) 与任意 \(y\ne h\) 有严格不等式
\[
\frac{(K^\star_\delta)^m(y\mid x)}{w(y)}>\frac{(K^\star_\delta)^m(h\mid x)}{w(h)}.
\]
\end{corollary}
\begin{proof}
若存在 \(y\ne h\) 使某个 \(m\) 处取等，则需有
\(\PP_x(X_m=y,\,T>m)=0\)。
由于 \(T\) 仅在刷新步停止，\(\PP_x(X_m=y,\,T>m)=0\) 蕴含：一旦在某次刷新采样到 \(y\) 即必然接受。
这要求 \(a_\delta(y)=1\)，即 \(r_\delta(y)=c_\delta\)，与 \(\mathcal H_\delta=\{h\}\) 矛盾。证毕。
\end{proof}

\begin{theorem}[强平稳时刻的母函数闭式]\label{thm:pom-diagonal-rate-accept-refresh-sst-pgf}
\leanverified{paper\_pom\_diagonal\_rate\_accept\_refresh\_sst\_pgf}
在定义 \ref{def:pom-diagonal-rate-accept-refresh-sst} 的设定下，
取 \(x\in X\setminus\{h\}\)。
定义
\[
\mathcal P_\delta(z):=\prod_{u\in X}(t_u-z),
\qquad
\mathcal Q_\delta(z):=z\,\mathcal P_\delta'(z)+\kappa_\delta\,\mathcal P_\delta(z),
\qquad
\mathcal P_{-x,-h}(z):=\prod_{u\in X\setminus\{x,h\}}(t_u-z).
\]
则对 \(|s|\) 充分小，有严格恒等式
\[
\EE_x[s^T]
=\frac{\kappa_\delta(1+\kappa_\delta)}{1+\kappa_\delta\delta}\,
(t_x-\kappa_\delta)(t_h-\kappa_\delta)\cdot
\frac{s(1-s)\,\mathcal P_{-x,-h}(\kappa_\delta s)}{\mathcal Q_\delta(\kappa_\delta s)}.
\]
\end{theorem}
\begin{proof}
以“刷新步”为唯一停止机会，\(\EE_x[s^T]\) 满足一组线性母函数方程 \((I-s\widetilde K)f=sq\)，其中 \(\widetilde K\) 为“未停止”下的次转移核。
在最优核参数化中，\(\widetilde K\) 保持“对角 \(+\) 秩一”结构，
从而可用秩一行列式引理把解压缩到一维多项式比值；
再利用恒等式 \(A_\delta^2=(1+\kappa_\delta\delta)^{-1}(1+\kappa_\delta)\)（见定理 \ref{thm:pom-diagonal-rate-diagonal-statistics-complete}）与
\(\mathcal Q_\delta(\kappa_\delta)=0\)（定理 \ref{thm:pom-maxent-markov-laguerre-secular-spectrum}）
化简即可得到闭式。证毕。
\end{proof}

\begin{corollary}[分离距离母函数的有理性]\label{cor:pom-diagonal-rate-separation-pgf-rational}
\leanverified{paper\_pom\_diagonal\_rate\_separation\_pgf\_rational}
在定理 \ref{thm:pom-diagonal-rate-accept-refresh-sst-pgf} 的设定下，
令
\[
\Delta_\delta(s):=\frac{\mathcal Q_\delta(\kappa_\delta s)}{1-s},
\qquad
\mathcal S_x(s):=\sum_{m\ge 0}\mathrm{sep}_x(m)\,s^m=\sum_{m\ge 0}\PP_x(T>m)\,s^m.
\]
则 \(\Delta_\delta\) 为次数 \(n-1\) 多项式，并且存在次数至多 \(n-2\) 的多项式 \(N_{x,\delta}\) 使
\[
\mathcal S_x(s)=\frac{N_{x,\delta}(s)}{\Delta_\delta(s)}.
\]
更精确地，可取
\[
N_{x,\delta}(s)
:=\frac{\Delta_\delta(s)
-\frac{\kappa_\delta(1+\kappa_\delta)}{1+\kappa_\delta\delta}(t_x-\kappa_\delta)(t_h-\kappa_\delta)\,
s\,\mathcal P_{-x,-h}(\kappa_\delta s)}{1-s}.
\]
\end{corollary}
\begin{proof}
由定理 \ref{thm:pom-diagonal-rate-accept-refresh-separation-exact} 与恒等式
\(
\sum_{m\ge 0}\PP(T>m)s^m=\frac{1-\EE[s^T]}{1-s}
\)
得
\(
\mathcal S_x(s)=\frac{1-\EE_x[s^T]}{1-s}
\)。
将定理 \ref{thm:pom-diagonal-rate-accept-refresh-sst-pgf} 代入，并用
\(\mathcal Q_\delta(\kappa_\delta s)=(1-s)\Delta_\delta(s)\)
整理即得。
由于 \(\mathcal S_x(1)=\EE_x[T]<\infty\)，分子在 \(s=1\) 处必消去，从而 \(N_{x,\delta}\) 为多项式且 \(\deg N_{x,\delta}\le n-2\)。证毕。
\end{proof}

\begin{theorem}[分离距离的谱展开与残数系数]\label{thm:pom-diagonal-rate-separation-spectral-residue}
\leanverified{paper\_pom\_diagonal\_rate\_separation\_spectral\_residue}
设 \(t_u\) 两两不同，并令 \(\mathcal Q_\delta\) 的全部零点按升序记为
\[
z_1=\kappa_\delta<z_2<\cdots<z_n.
\]
置
\(
\lambda_i:=\kappa_\delta/z_i\in(0,1)
\)
（\(i=2,\dots,n\)）。
则对任意 \(x\ne h\) 与任意 \(m\ge 0\)，存在唯一系数 \(\{\omega_i(x)\}_{i=2}^{n}\subset\RR\)（\(\sum_{i=2}^n\omega_i(x)=1\)）使
\[
\mathrm{sep}_x(m)=\sum_{i=2}^{n}\omega_i(x)\,\lambda_i^{m}.
\]
并且系数具有显式残数闭式
\[
\omega_i(x)=\frac{1+\kappa_\delta}{1+\kappa_\delta\delta}\,
(t_x-\kappa_\delta)(t_h-\kappa_\delta)\cdot
\frac{\mathcal P_{-x,-h}(z_i)}{\mathcal Q_\delta'(z_i)}
\qquad(i=2,\dots,n).
\]
\end{theorem}
\begin{proof}
由推论 \ref{cor:pom-diagonal-rate-separation-pgf-rational}，
\(\mathcal S_x(s)=N_{x,\delta}(s)/\Delta_\delta(s)\) 为有理函数，
其极点为 \(\Delta_\delta(s)=0\) 的根 \(s_i=z_i/\kappa_\delta=1/\lambda_i\)（\(i=2,\dots,n\)），且均为单极点。
对 \(\mathcal S_x\) 作部分分式分解得到
\[
\mathcal S_x(s)=\sum_{i=2}^{n}\frac{\omega_i(x)}{1-\lambda_i s},
\]
从而系数展开给出 \(\mathrm{sep}_x(m)=\sum_{i=2}^{n}\omega_i(x)\lambda_i^m\)；
唯一性由 \(\{\lambda_i\}\) 互异的线性无关性给出。
\par
残数闭式由
\(
\mathrm{Res}_{s=s_i}\mathcal S_x(s)=-\omega_i(x)/\lambda_i
\)
及
\(
\Delta_\delta'(s_i)=\kappa_\delta\,\mathcal Q_\delta'(z_i)/(1-s_i)
\)
与 \(N_{x,\delta}(s_i)=-C_\delta\,s_i\,\mathcal P_{-x,-h}(z_i)/(1-s_i)\)
（其中 \(C_\delta=\kappa_\delta(1+\kappa_\delta)(t_x-\kappa_\delta)(t_h-\kappa_\delta)/(1+\kappa_\delta\delta)\)）
代回并化简得到。证毕。
\end{proof}

\begin{proposition}[系数符号的交错判据]\label{prop:pom-diagonal-rate-separation-weight-sign-rule}
\leanverified{paper\_pom\_diagonal\_rate\_separation\_weight\_sign\_rule}
在定理 \ref{thm:pom-diagonal-rate-separation-spectral-residue} 的设定下，
对每个 \(i\ge 2\) 有符号恒等式
\[
\mathrm{sign}\bigl(\omega_i(x)\bigr)=\mathrm{sign}(t_x-z_i).
\]
特别地，若 \(x\in\arg\max_{u\in X}t_u\)（等价地 \(p_\delta(x)=K^\star_\delta(x\mid x)\) 取最小值），则 \(\omega_i(x)>0\) 对所有 \(i=2,\dots,n\) 成立。
\end{proposition}
\begin{proof}
由定理 \ref{thm:pom-maxent-markov-laguerre-secular-spectrum} 的交错性，\(z_i\notin\{t_u\}\) 且 \(z_i\) 为单根。
记
\(
\phi(z):=\sum_{u}(t_u-z)^{-1}-\kappa_\delta/z
\)。
则 \(\phi(z)=0\) 与 \(\mathcal Q_\delta(z)=0\) 等价，并且 \(\phi'(z)=\sum_u(t_u-z)^{-2}+\kappa_\delta/z^2>0\)。
利用恒等式 \(\phi(z)=-\mathcal Q_\delta(z)/(z\mathcal P_\delta(z))\)，在零点处得到
\[
\mathcal Q_\delta'(z_i)=-z_i\,\mathcal P_\delta(z_i)\,\phi'(z_i),
\]
从而 \(\mathrm{sign}(\mathcal Q_\delta'(z_i))=-\mathrm{sign}(\mathcal P_\delta(z_i))\)。
另一方面，
\(
\mathcal P_\delta(z_i)=(t_x-z_i)(t_h-z_i)\mathcal P_{-x,-h}(z_i)
\)。
由于 \(h\in\arg\min t\) 且 \(i\ge 2\) 蕴含 \(z_i>t_h\)，有 \(t_h-z_i<0\)。
将上述符号关系代入定理 \ref{thm:pom-diagonal-rate-separation-spectral-residue} 的残数闭式即得
\(\mathrm{sign}(\omega_i(x))=\mathrm{sign}(t_x-z_i)\)。证毕。
\end{proof}

\begin{corollary}[极端起点下的完全单调与严格对数凸]\label{cor:pom-diagonal-rate-separation-hausdorff-logconvex}
\leanverified{paper\_pom\_diagonal\_rate\_separation\_hausdorff\_logconvex}
在定理 \ref{thm:pom-diagonal-rate-separation-spectral-residue} 的设定下，
若 \(x\in\arg\max_{u}t_u\)，则 \(\{\mathrm{sep}_x(m)\}_{m\ge 0}\) 为 \([0,1)\) 上离散概率测度
\(
\sum_{i=2}^{n}\omega_i(x)\delta_{\lambda_i}
\)
的 Hausdorff 矩序列，因而对任意 \(k,m\ge 0\) 成立完全单调性
\[
(-1)^k(\Delta^k \mathrm{sep}_x)_m\ge 0.
\]
并且当 \(n\ge 3\) 时，对所有 \(m\ge 1\) 有严格对数凸与比值单调：
\[
\mathrm{sep}_x(m)^2<\mathrm{sep}_x(m-1)\mathrm{sep}_x(m+1),
\qquad
\frac{\mathrm{sep}_x(m+1)}{\mathrm{sep}_x(m)}\ \text{\rm 严格递增}.
\]
\end{corollary}
\begin{proof}
由命题 \ref{prop:pom-diagonal-rate-separation-weight-sign-rule}，
\(\omega_i(x)>0\) 且 \(\sum_{i=2}^n\omega_i(x)=1\)，并由定理 \ref{thm:pom-maxent-markov-laguerre-secular-spectrum} 知 \(\lambda_i\in(0,1)\)。
故由定理 \ref{thm:pom-diagonal-rate-separation-spectral-residue} 的表示得到 Hausdorff 矩序列表述与完全单调性。
当 \(n\ge 3\) 时，至少存在两条不同的 \(\lambda_i\)，从而完全单调矩序列的标准严格性结论给出严格对数凸与比值单调。证毕。
\end{proof}

\leanverified{paper\_pom\_diagonal\_rate\_diagonal\_determines\_separation\_profile}
\begin{corollary}[单次对角观测确定分离距离剖面]\label{cor:pom-diagonal-rate-diagonal-determines-separation-profile}
在定理 \ref{thm:pom-diagonal-rate-diagonal-statistics-complete} 的设定下，
仅由对角保持率向量 \(p_\delta(x)=K^\star_\delta(x\mid x)\) 即可回收 \(\kappa_\delta\) 与 \(\{t_x\}_{x\in X}\)，
从而可构造 \(\mathcal P_\delta,\mathcal Q_\delta\) 并由定理 \ref{thm:pom-diagonal-rate-separation-spectral-residue} 给出任意起点 \(x\ne h\) 的精确分离距离序列 \(\{\mathrm{sep}_x(m)\}_{m\ge 0}\)。
\end{corollary}


\paragraph{删点基本矩阵的秩一闭式}
对删点链的占据与访问次数，秩一刷新结构同样给出一行公式封闭。

\begin{theorem}[删点基本矩阵秩一闭式]\label{thm:pom-diagonal-rate-refresh-fundamental-matrix-rankone}
\leanverified{paper\_pom\_diagonal\_rate\_refresh\_fundamental\_matrix\_rankone}
在秩一刷新表示 \(K_\delta=\mathrm{diag}(1-r_\delta)+r_\delta\,\pi_\delta^{\mathsf T}\) 下，固定 \(y\in X\)，
令 \(T^{(y)}\) 表示在 \(X\setminus\{y\}\) 上由删去 \(y\) 诱导的次转移矩阵，并定义基本矩阵 \(N^{(y)}:=(I-T^{(y)})^{-1}\)。
则对任意 \(x,z\in X\setminus\{y\}\) 有严格恒等式
\[
N^{(y)}_{xz}=\frac{1}{r_\delta(x)}\,\mathbf 1_{\{x=z\}}+\frac{\pi_\delta(z)}{\pi_\delta(y)\,r_\delta(z)}.
\]
从而对固定 \(z\ne y\)，期望访问次数 \(N^{(y)}_{xz}\) 对起点 \(x\) 的依赖仅通过初始补丁项 \(\frac{1}{r_\delta(x)}\mathbf 1_{\{x=z\}}\) 出现；其余部分与 \(x\) 无关。
并且对 \(x\ne y\) 有
\(
\EE_x[\tau_y]=\sum_{z\ne y}N^{(y)}_{xz}
\)，从而推出推论 \ref{cor:pom-diagonal-rate-refresh-hitting-time-mean-closed}。
\end{theorem}
\begin{proof}
在 \(X\setminus\{y\}\) 上，次转移矩阵可写为
\[
T^{(y)}=\mathrm{diag}(1-r_\delta)+r_\delta\,\pi_{-y}^{!\mathsf T},
\]
其中 \(\pi_{-y}^{!}\) 为删去 \(y\) 的行向量。
因此
\[
I-T^{(y)}=\mathrm{diag}(r_\delta)-r_\delta\,\pi_{-y}^{!\mathsf T}.
\]
记 \(D:=\mathrm{diag}(r_\delta)\)（可逆），并注意 \(\pi_{-y}^{!\mathsf T}D^{-1}r_\delta=\sum_{z\ne y}\pi_\delta(z)=1-\pi_\delta(y)\)。
对 \(D-r_\delta\pi_{-y}^{!\mathsf T}\) 应用 Sherman--Morrison 公式得到
\[
(I-T^{(y)})^{-1}
=D^{-1}+\frac{D^{-1}r_\delta\,\pi_{-y}^{!\mathsf T}D^{-1}}{1-(1-\pi_\delta(y))}.
\]
并用 \(D^{-1}r_\delta=\mathbf 1\) 化简即可得到所示矩阵元公式。
最后对 \(z\) 求和得到 \(\EE_x[\tau_y]=\sum_{z\ne y}N^{(y)}_{xz}\)，并与前式比较即得推论 \ref{cor:pom-diagonal-rate-refresh-hitting-time-mean-closed}。证毕。
\end{proof}

\paragraph{多靶集合的闭式推广}
上述删点结构对任意非空真子集 \(B\subset X\) 保持同型：删集 Laguerre 多项式同时控制命中时分布与基本矩阵。

\begin{theorem}[删集 Laguerre 多项式控制 \texorpdfstring{$\tau_B$}{tau_B}]\label{thm:pom-diagonal-rate-refresh-hitting-time-deleted-set-laguerre-pgf}
在命题 \ref{prop:pom-diagonal-rate-diagonal-closure-rankone} 的设定下，令 \(B\subset X\) 为非空真子集，定义命中时
\(
\tau_B:=\inf\{t\ge 0:\ X_t\in B\}
\)。
记
\[
P_{-B}(z):=\prod_{u\in X\setminus B}(t_u-z),\qquad
D_B(z):=zP_{-B}'(z)+\kappa_\delta P_{-B}(z),
\qquad
\pi_\delta(B):=\sum_{b\in B}\pi_\delta(b)=\sum_{b\in B}\frac{1}{t_b-\kappa_\delta}.
\]
则对任意 \(x\notin B\) 与 \(|s|\) 充分小，有严格恒等式
\[
\EE_x[s^{\tau_B}]
=
\kappa_\delta\,s\,(t_x-\kappa_\delta)\,\pi_\delta(B)\cdot
\frac{P_{-(B\cup\{x\})}(\kappa_\delta s)}{D_B(\kappa_\delta s)}.
\]
\end{theorem}
\begin{proof}
将定理 \ref{thm:pom-diagonal-rate-refresh-hitting-time-pgf-closed} 的论证对目标集合 \(B\) 平行改写即可：
失败轮数服从 \(\mathrm{Geom}(\pi_\delta(B))\)，且失败落点为条件化分布 \(\pi_\delta(\cdot\mid \cdot\notin B)\) 的独立样本。
在最优核参数化下，
\(
\sum_{z\notin B}\pi_\delta(z)G_z(s)=s\sum_{z\notin B}(t_z-\kappa_\delta s)^{-1}
\)；
令 \(z=\kappa_\delta s\) 并用对数导数恒等式
\(
\sum_{u\notin B}(t_u-z)^{-1}=-P_{-B}'(z)/P_{-B}(z)
\)
化简得到
\(
\bigl(1-s\sum_{z\notin B}(t_z-\kappa_\delta s)^{-1}\bigr)^{-1}=\kappa_\delta P_{-B}(\kappa_\delta s)/D_B(\kappa_\delta s)
\)。
再用 \(P_{-B}(\kappa_\delta s)=(t_x-\kappa_\delta s)P_{-(B\cup\{x\})}(\kappa_\delta s)\) 消去 \((t_x-\kappa_\delta s)^{-1}\) 即得。证毕。
\end{proof}
\leanverified{paper\_pom\_diagonal\_rate\_refresh\_hitting\_time\_deleted\_set\_laguerre\_pgf}

\begin{proposition}[删集基本矩阵的秩一闭式]\label{prop:pom-diagonal-rate-refresh-fundamental-matrix-rankone-set}
\leanverified{paper\_pom\_diagonal\_rate\_refresh\_fundamental\_matrix\_rankone\_set}
在定理 \ref{thm:pom-diagonal-rate-refresh-fundamental-matrix-rankone} 的设定下，令 \(B\subset X\) 为非空真子集，记 \(T^{(B)}\) 为在 \(X\setminus B\) 上删去 \(B\) 后诱导的次转移矩阵，\(N^{(B)}:=(I-T^{(B)})^{-1}\) 为基本矩阵。
则对任意 \(x,z\notin B\) 有
\[
N^{(B)}_{xz}=\frac{1}{r_\delta(x)}\,\mathbf 1_{\{x=z\}}+\frac{\pi_\delta(z)}{\pi_\delta(B)\,r_\delta(z)}.
\]
\end{proposition}
\begin{proof}
与定理 \ref{thm:pom-diagonal-rate-refresh-fundamental-matrix-rankone} 同型：在删集情形下仅需将 \(\pi_\delta(y)\) 替换为 \(\pi_\delta(B)\)，并将求和索引从 \(X\setminus\{y\}\) 改为 \(X\setminus B\)；
Sherman--Morrison 化简不变。证毕。
\end{proof}

\begin{conjecture}[Laguerre--P\'olya 闭包与极限实零性]\label{conj:pom-diagonal-rate-refresh-deleted-laguerre-lp-closure}
设存在一族规模趋于无穷的有限模型
\(
(X_m,\kappa_m,\{t_x^{(m)}\}_{x\in X_m})
\)
及靶集合 \(B_m\subset X_m\)（非空真子集），使删集多项式
\[
D_{B_m}^{(m)}(z):=z\,(P_{-B_m}^{(m)})'(z)+\kappa_m P_{-B_m}^{(m)}(z),
\qquad
P_{-B_m}^{(m)}(z):=\prod_{u\in X_m\setminus B_m}(t_u^{(m)}-z)
\]
在适当尺度归一化后局部一致收敛到整函数 \(D_\infty(z)\)。
若经验谱测度
\[
\nu_m:=\frac{1}{|X_m\setminus B_m|}\sum_{x\notin B_m}\delta_{t_x^{(m)}}
\]
在某个 \((\underline t,\infty)\) 上弱收敛到极限测度 \(\nu\)，并满足一条 Stieltjes 紧性条件（例如 \(\sup_m\int (t-\underline t)^{-2}\,d\nu_m(t)<\infty\) 且 \(\inf_m\kappa_m>0\)），
则极限整函数 \(D_\infty\) 属于 Laguerre--P\'olya 类，
其全部零点为正实数并与 \(\mathrm{supp}(\nu)\) 的间隙保持交错。
该实零性质在删集命中时的几何谱表示中对应于谱参数落在 \((0,1)\)；
而系数测度的正性（从而 Hausdorff 尾矩与完全单调性）则属于独立的附加约束，需在具体折叠/极限构造中另行核验。
\end{conjecture}

\begin{proposition}[有向命中时的秩二分解与势差闭合]\label{prop:pom-diagonal-rate-refresh-hitting-time-rank-two-potential}
\leanverified{paper\_pom\_diagonal\_rate\_refresh\_hitting\_time\_rank\_two\_potential}
对互异 \(x,y\in X\) 定义
\(
m_\delta(x,y):=\EE_x[\tau_y]
\)。
则存在函数 \(A_\delta:X\to(0,\infty)\) 与 \(B_\delta:X\to(0,\infty)\)，使得对一切 \(x\ne y\) 有分离表示
\[
m_\delta(x,y)=A_\delta(x)+B_\delta(y),
\qquad
A_\delta(x)=\frac{1}{r_\delta(x)},
\quad
B_\delta(y)=\frac{1}{\pi_\delta(y)}\sum_{z\ne y}\frac{\pi_\delta(z)}{r_\delta(z)}.
\]
从而非对角均值命中时矩阵 \(\bigl(m_\delta(x,y)\bigr)_{x\ne y}\) 的秩至多为 \(2\)。
并且定义势函数 \(\Phi_\delta(x):=A_\delta(x)-B_\delta(x)\)，则对任意互异 \(x,y\in X\) 有严格恒等式
\[
m_\delta(x,y)-m_\delta(y,x)=\Phi_\delta(x)-\Phi_\delta(y),
\]
从而对任意有向环 \(x_1\to x_2\to\cdots\to x_k\to x_1\) 成立环路守恒
\[
\sum_{i=1}^{k} m_\delta(x_i,x_{i+1})=\sum_{i=1}^{k} m_\delta(x_{i+1},x_i),
\qquad (x_{k+1}:=x_1).
\]
\end{proposition}
\begin{proof}
由推论 \ref{cor:pom-diagonal-rate-refresh-hitting-time-mean-closed} 逐点分离即得 \(m_\delta(x,y)=A_\delta(x)+B_\delta(y)\)。
势差恒等式由代入并消去对称项立即推出；环路守恒为沿环求和后的直接推论。证毕。
\end{proof}

\paragraph{电导完全图的星形 Kron 还原与通勤时闭式}
由最优耦合的详细平衡，核 \(K_\delta\) 在平稳分布 \(w\) 下可逆。
定义对称电导（边权）
\[
c_{xy}:=w(x)K_\delta(y\mid x)\qquad(x\ne y),
\]
并令
\[
S_\delta:=\sum_{z\in X}\frac{\pi_\delta(z)}{r_\delta(z)}\ (>0).
\]

\begin{proposition}[秩一电导与星形网络等价]\label{prop:pom-diagonal-rate-refresh-kron-star-equivalence}
\leanverified{paper\_pom\_diagonal\_rate\_refresh\_kron\_star\_equivalence}
在上述记号下，对任意互异 \(x,y\in X\) 有严格恒等式
\[
c_{xy}=\frac{1}{S_\delta}\,\pi_\delta(x)\pi_\delta(y).
\]
从而完全图电导为秩一乘积。
进一步，考虑星形图 \(S\)（枢纽 \(0\) 与叶集 \(X\)），对边 \((0,x)\) 赋电导
\[
g_x:=\frac{\pi_\delta(x)}{S_\delta}.
\]
则消去枢纽 \(0\) 的 Kron 还原（Schur 补）在叶集 \(X\) 上诱导的等效完全图电导恰为 \(\{c_{xy}\}_{x\ne y}\)。
\end{proposition}
\begin{proof}
由命题 \ref{prop:pom-diagonal-rate-rank-one-refresh-decomposition}，对 \(x\ne y\) 有 \(K_\delta(y\mid x)=r_\delta(x)\pi_\delta(y)\)，故
\[
c_{xy}=w(x)r_\delta(x)\pi_\delta(y).
\]
对 \(x\ne y\) 的详细平衡 \(w(x)K_\delta(y\mid x)=w(y)K_\delta(x\mid y)\) 给出
\(
w(x)r_\delta(x)/\pi_\delta(x)=w(y)r_\delta(y)/\pi_\delta(y)
\)，
从而存在常数 \(C_\delta>0\) 使 \(w(x)r_\delta(x)=C_\delta\,\pi_\delta(x)\)。
对 \(x\) 求和并用 \(\sum_x w(x)=1\) 得
\[
1=\sum_x w(x)=\sum_x \frac{C_\delta\,\pi_\delta(x)}{r_\delta(x)}=C_\delta\,S_\delta,
\qquad\text{即}\qquad C_\delta=\frac{1}{S_\delta}.
\]
代回得 \(c_{xy}=S_\delta^{-1}\pi_\delta(x)\pi_\delta(y)\)。
\par
对星形图的 Kron 还原，标准 Schur 补给出叶间等效电导
\(
c^{\mathrm{eff}}_{xy}=g_xg_y/\sum_{z\in X}g_z
\)。
由于 \(\sum_z g_z=S_\delta^{-1}\sum_z\pi_\delta(z)=S_\delta^{-1}\)，得
\(
c^{\mathrm{eff}}_{xy}=S_\delta^{-1}\pi_\delta(x)\pi_\delta(y)=c_{xy}
\)，即为所示。证毕。
\end{proof}

\begin{corollary}[有效电阻与通勤时的星形闭式]\label{cor:pom-diagonal-rate-refresh-commute-effective-resistance}
记 \(R_{\mathrm{eff}}(x,y)\) 为上述电网络在叶集 \(X\) 上的有效电阻，则对任意互异 \(x,y\in X\) 有
\[
R_{\mathrm{eff}}(x,y)=\frac{1}{g_x}+\frac{1}{g_y}
=S_\delta\left(\frac{1}{\pi_\delta(x)}+\frac{1}{\pi_\delta(y)}\right).
\]
并且通勤时
\(
\mathrm{Comm}(x,y):=\EE_x[\tau_y]+\EE_y[\tau_x]
\)
满足严格闭式
\[
\mathrm{Comm}(x,y)=S_\delta\left(\frac{1}{\pi_\delta(x)}+\frac{1}{\pi_\delta(y)}\right)
=R_{\mathrm{eff}}(x,y).
\]
\end{corollary}
\begin{proof}
有效电阻恒等式由星形网络两叶间串联电阻 \(1/g_x+1/g_y\) 直接得到。
另一方面，由推论 \ref{cor:pom-diagonal-rate-refresh-hitting-time-mean-closed} 得
\[
\EE_x[\tau_y]+\EE_y[\tau_x]
=\frac{1}{\pi_\delta(y)}\sum_{z\ne y}\frac{\pi_\delta(z)}{r_\delta(z)}
\ +\ \frac{1}{\pi_\delta(x)}\sum_{z\ne x}\frac{\pi_\delta(z)}{r_\delta(z)}.
\]
将 \(\sum_{z\ne y}\pi_\delta(z)/r_\delta(z)=S_\delta-\pi_\delta(y)/r_\delta(y)\) 代入并化简，得到
\(
\mathrm{Comm}(x,y)=S_\delta(\pi_\delta(x)^{-1}+\pi_\delta(y)^{-1})
\)。
与前式比较即得 \(\mathrm{Comm}(x,y)=R_{\mathrm{eff}}(x,y)\)。证毕。
\end{proof}

\paragraph{加权生成树总权重与伪行列式闭式}
令 \(n:=|X|\ge 2\)，并令 \(\tau_\delta\) 表示由 \(\{c_{xy}\}_{x\ne y}\) 定义的加权生成树总权重：
\[
\tau_\delta:=\sum_{T\ \text{\rm 生成树}}\ \prod_{\{x,y\}\in E(T)} c_{xy}.
\]

\begin{theorem}[加权 Cayley--Pr\"ufer 闭式]\label{thm:pom-diagonal-rate-refresh-weighted-cayley-prufer}
在命题 \ref{prop:pom-diagonal-rate-refresh-kron-star-equivalence} 的设定下，有严格恒等式
\[
\tau_\delta=S_\delta^{-(n-1)}\prod_{x\in X}\pi_\delta(x).
\]
\end{theorem}
\begin{proof}
由命题 \ref{prop:pom-diagonal-rate-refresh-kron-star-equivalence}，对 \(x\ne y\) 有
\(
c_{xy}=a_xa_y
\)
其中 \(a_x:=\pi_\delta(x)/\sqrt{S_\delta}\)。
对边权 \(a_xa_y\) 的完全图，Pr\"ufer 码计数给出加权 Cayley 恒等式
\(
\sum_T\prod_{\{x,y\}\in E(T)}a_xa_y=(\sum_x a_x)^{n-2}\prod_x a_x
\)。
代入 \(\sum_x a_x=1/\sqrt{S_\delta}\) 即得所示闭式。证毕。
\end{proof}

\begin{corollary}[归一化 Laplacian 的伪行列式]\label{cor:pom-diagonal-rate-refresh-laplacian-pdet}
令 \(L\) 为该完全图的 Laplacian（对 \(x\ne y\) 取 \(L_{xy}=-c_{xy}\)，并令 \(L_{xx}=\sum_{y\ne x}c_{xy}\)）。
则矩阵树定理给出
\[
\operatorname{pdet}(L)=n\,\tau_\delta.
\]
进一步，对归一化 Laplacian \(\mathcal L:=\mathrm{diag}(w)^{-1/2}L\,\mathrm{diag}(w)^{-1/2}\) 有
\[
\operatorname{pdet}(\mathcal L)=\frac{\tau_\delta}{\prod_{x\in X}w(x)}.
\]
\end{corollary}
\begin{proof}
第一式为矩阵树定理的标准形式：\(\det(L^{(x)})=\tau_\delta\)（删去第 \(x\) 行列的任一主余子式），并由 \(L\) 的零空间由 \(\mathbf 1\) 张成而得 \(\operatorname{pdet}(L)=n\tau_\delta\)。
\par
对第二式，注意 \(\mathcal L\) 的零向量为 \(\sqrt w\)（逐点平方根），且任取 \(x\) 的主余子式满足
\[
\det(\mathcal L^{(x)})=\frac{\det(L^{(x)})}{\prod_{z\ne x}w(z)}=\frac{\tau_\delta}{\prod_{z\ne x}w(z)}.
\]
由秩 \(n-1\) 对称矩阵的伴随矩阵结构，\(\det(\mathcal L^{(x)})/\bigl(\sqrt{w(x)}\,\bigr)^2\) 与 \(x\) 无关，并等于 \(\operatorname{pdet}(\mathcal L)/\norm{\sqrt w}_2^2\)；
而 \(\norm{\sqrt w}_2^2=\sum_x w(x)=1\)，故
\(
\operatorname{pdet}(\mathcal L)=\tau_\delta/\prod_x w(x)
\)。
证毕。
\end{proof}

\paragraph{单次对角观测决定通勤时与 Kemeny 常数}
对角保持率 \(p_\delta(x)=K^\star_\delta(x\mid x)\) 的完备性已由定理 \ref{thm:pom-diagonal-rate-diagonal-statistics-complete} 给出。
结合推论 \ref{cor:pom-diagonal-rate-refresh-commute-effective-resistance} 可将通勤时矩阵与 Kemeny 常数闭合为单次对角统计的代数函数。

\begin{corollary}[对角保持率决定全通勤时矩阵]\label{cor:pom-diagonal-rate-diagonal-determines-commute}
沿用定理 \ref{thm:pom-diagonal-rate-diagonal-statistics-complete} 的记号，
由对角保持率向量 \(p_\delta(\cdot)\) 恢复 \(\kappa_\delta\) 与 \(\pi_\delta\)，并令
\[
r_\delta(x):=1-\frac{\kappa_\delta}{1+\kappa_\delta}\,p_\delta(x),
\qquad
S_\delta:=\sum_{z\in X}\frac{\pi_\delta(z)}{r_\delta(z)}.
\]
则对任意互异 \(x,y\in X\) 有通勤时闭式
\[
\mathrm{Comm}(x,y)=S_\delta\left(\frac{1}{\pi_\delta(x)}+\frac{1}{\pi_\delta(y)}\right).
\]
\end{corollary}
\begin{proof}
由定理 \ref{thm:pom-diagonal-rate-diagonal-statistics-complete}，\(\kappa_\delta,\pi_\delta\) 与 \(K^\star_\delta\) 可由 \(p_\delta(\cdot)\) 显式回收；
并由 \(K^\star_\delta=\mathrm{diag}(1-r_\delta)+r_\delta\pi_\delta^{\mathsf T}\) 得 \(r_\delta(x)=1-\frac{\kappa_\delta}{1+\kappa_\delta}p_\delta(x)\)。
代入推论 \ref{cor:pom-diagonal-rate-refresh-commute-effective-resistance} 即得。证毕。
\end{proof}

\begin{theorem}[Kemeny 常数的星形闭式]\label{thm:pom-diagonal-rate-refresh-kemeny-star-closed}
\leanverified{paper\_pom\_diagonal\_rate\_refresh\_kemeny\_star\_closed}
令
\(
\mathcal K_\delta:=\sum_{y\in X}w(y)\,\EE_x[\tau_y]
\)
为 Kemeny 常数（其与起点 \(x\) 无关）。
则在上述记号下成立严格恒等式
\[
\mathcal K_\delta
=S_\delta\sum_{y\in X}\frac{w(y)\bigl(1-w(y)\bigr)}{\pi_\delta(y)}.
\]
并且借助定理 \ref{thm:pom-diagonal-rate-diagonal-statistics-complete} 的对角闭包，\(\mathcal K_\delta\) 可由单次对角观测 \(p_\delta(\cdot)\) 在 \(O(n)\) 时间内精确回收。
\end{theorem}
\begin{proof}
由命题 \ref{prop:pom-diagonal-rate-refresh-hitting-time-rank-two-potential} 的秩二分解，
对任意起点 \(x\) 有
\[
\sum_{y\in X}w(y)\,\EE_x[\tau_y]
=\sum_{y\ne x}w(y)\bigl(A_\delta(x)+B_\delta(y)\bigr)
=A_\delta(x)\bigl(1-w(x)\bigr)+\sum_{y\ne x}w(y)B_\delta(y).
\]
利用 \(w(x)r_\delta(x)=S_\delta^{-1}\pi_\delta(x)\)（命题 \ref{prop:pom-diagonal-rate-refresh-kron-star-equivalence} 的证明）可验证右端与 \(x\) 无关，从而该和确为 Kemeny 常数。
另一方面，将
\(
B_\delta(y)=\pi_\delta(y)^{-1}\sum_{z\ne y}\pi_\delta(z)/r_\delta(z)
=S_\delta/\pi_\delta(y)-1/r_\delta(y)
\)
代回并整理得
\[
\mathcal K_\delta
=S_\delta\sum_{y\ne x}\frac{w(y)}{\pi_\delta(y)}-\sum_{y\ne x}\frac{w(y)}{r_\delta(y)}
=S_\delta\sum_{y\in X}\frac{w(y)\bigl(1-w(y)\bigr)}{\pi_\delta(y)},
\]
其中最后一步使用恒等式
\(
\sum_y w(y)/r_\delta(y)=S_\delta\sum_y w(y)^2/\pi_\delta(y)
\)
（由 \(1/r_\delta(y)=S_\delta w(y)/\pi_\delta(y)\) 得出）。
时间复杂度陈述由 \(\kappa_\delta,\pi_\delta,w\) 的对角闭式与一次线性求和直接推出。证毕。
\par
可复现核验脚本：\texttt{scripts/exp\_pom\_diagonal\_rate\_hitting\_star\_commute\_audit.py}。
\end{proof}
